% ---------------------------------------------------------------------------
% Section: Training through time: the chained expansion, its applications,
%          and why the late stage is not predictable from the initialisation
%
% Drop-in section for the paper. Conventions of Part III: Pi = I, c_p = 1,
% gradient flow, E = \hat Y - Y, K(theta) the (1/n-normalised) tangent kernel
% so that \dot E = -K E, J the Jacobian on the data, theta_0 the initialisation,
% lambda_1 = lambda_max(K), H_2 the second-derivative bound of Part I.
% Cross-references to earlier results (Proposition 20 = the analyticity radius
% of expansions about theta_0; Section sec:held-rough = the held rough error)
% are hard-coded; replace them with \ref{...} to your labels. Environments:
% theorem, proposition, corollary, definition, assumption, remark, proof.
% Every numerical statement below was obtained at n = 200 training points on
% the Lorenz-63 testbed and is marked [n=200]; see the last subsection.
% ---------------------------------------------------------------------------

\section{Training through time: the chained expansion and the limits of prediction}
\label{sec:chain}

Proposition~20 says that no series in $\theta-\theta_0$ reaches the stretch of
training in which the rough part of the error is removed, and
Section~\ref{sec:held-rough} replaced the expansion variable by the peak
feature roughness. This section keeps $\theta$ as the variable and moves the
expansion point instead: training is written as a chain of expansions, each
about the current state and each valid within its own radius, whose
composition is exact. The chain gives three kinds of statement. Identities
that hold at every time for any network in a broad class (\S\ref{ssec:chain-identities});
structure theorems for what gradient flow does to the weights and the kernel
at each stage, which is where the low-rank ``spikes'' and the targeted growth
of the kernel come from (\S\ref{ssec:chain-spikes}--\ref{ssec:chain-richness});
and applications in which an open question becomes a computation with the
chain or a theorem about its leading orders (\S\ref{ssec:chain-apps}). The
last two subsections state why the late stage cannot be predicted from the
initialisation without carrying the state forward, in a form that is a
theorem rather than an impression (\S\ref{ssec:chain-impossible}), and what a
statistical description can and cannot promise (\S\ref{ssec:chain-statistical}).

\subsection{Setting and exact identities}
\label{ssec:chain-identities}

\paragraph{Model class.} $N_\theta$ is any composition of \emph{parameter-linear
blocks} $z_\ell = W_\ell a_\ell$, where $a_\ell$ is whatever the block receives
(dense, convolution, attention projections, normalisation gains), and
\emph{parameter-free maps} between them (pointwise activations, softmax,
normalisation), differentiable almost everywhere in $\theta$. ReLU qualifies.
Training is gradient flow on the squared loss, or any update of the
data-indexed outer-product form below (GD, SGD on a subset of the data,
momentum through its integral); Adam is treated in \S\ref{ssec:chain-apps}.
Write $\Psi(t)=J(\theta(t))\in\mathbb R^{N\times P}$ for the Jacobian on the $n$
data points ($N=n\,n_y$), $K=\Psi\Psi^{\!\top}/n$, $\delta_\ell^c(x)=\partial
\hat Y_c(x)/\partial z_\ell$ for the backprop signal into block $\ell$, and
$U(t)=\int_0^t E$ for the integrated residual.

\begin{proposition}[Four identities]
\label{prop:chain-identities}
For every $t$ at which $N_\theta$ is differentiable:
\begin{enumerate}
\item[(a)] $\dot E=-K(t)E$; hence $\tfrac{d}{dt}\|E\|^2=-2\langle E,KE\rangle$
and $E(T)=\mathrm{T}\!\exp\!\big(-\!\int_0^T K\big)E(0)$.
\item[(b)] $K=\frac1n\sum_\ell A_\ell\odot D_\ell$, with $A_\ell[i,j]=a_\ell(x_i)\!\cdot\!a_\ell(x_j)+1$
the Gram of the block's inputs and
$D_\ell[(i,c),(j,d)]=\sum_q\delta^c_{\ell,q}(x_i)\,\delta^d_{\ell,q}(x_j)$ the Gram
of its backprop signals.
\item[(c)] $W_\ell(t)-W_\ell(0)=-\frac1n\int_0^t\sum_i\sum_c\delta_\ell^c(x_i)\,E_c(x_i)\,a_\ell(x_i)^{\!\top}$:
the change of every block is a sum of outer products indexed by the data, its
row space lies in the span of the block's inputs on the data, and every block is
driven by the same residual.
\item[(d)] For a unit $q$ of block $\ell$ with preactivation $z_q$ and incoming
backward signal $m_q=W_{\ell+1}^{\!\top}\delta_{\ell+1}$,
\[
\dot z_q=\dot W_\ell a_\ell+W_\ell\,\dot a_\ell,\qquad
\dot m_q=\dot W_{\ell+1}^{\!\top}\delta_{\ell+1}+W_{\ell+1}^{\!\top}\dot\delta_{\ell+1}.
\]
The first terms are given by (c) (the \emph{coherent} part); the second carry
the motion of the neighbouring blocks through the current weights (the
\emph{propagation} part). Together with (a) and (b) this is a closed system on
the data.
\end{enumerate}
\end{proposition}
\begin{proof}
(a) is the chain rule for $\theta'=-\frac1nJ^{\!\top}E$. (b): $\partial\hat
Y_c(x)/\partial W_\ell[q,p]=\delta^c_{\ell,q}(x)\,a_{\ell,p}(x)$. (c): integrate
the gradient. (d): differentiate the definitions.
\end{proof}

\begin{remark}
Everything in this section is a consequence of
Proposition~\ref{prop:chain-identities}. In particular the propagation part of
(d) is the only place where the initialisation enters other than through the
unit statistics; omitting it (as any model built from the kernel alone must)
was measured to cost $11$--$28\%$ of the unit displacements early and a factor
$2$ in the late error on the reference problem [n=200].
\end{remark}

\subsection{The stage identity and the chain}
\label{ssec:chain-stage}

Fix stage boundaries $0=t_0<t_1<\dots<t_S=T$ and write
$\delta\Psi_s=\Psi(t_s)-\Psi(t_{s-1})$.

\begin{proposition}[Stage identity]
\label{prop:stage}
$K(t_s)-K(t_{s-1})=\mathrm{Cross}_s+Q_s$ with
$Q_s=\delta\Psi_s\delta\Psi_s^{\!\top}/n\succeq0$ and
$\mathrm{Cross}_s=(\Psi\delta\Psi_s^{\!\top}+\delta\Psi_s\Psi^{\!\top})/n$,
$\Psi=\Psi(t_{s-1})$. For a unit vector $v$, the quadratic form changes by
$\frac1n\big[2\langle\Psi^{\!\top}v,\delta\Psi_s^{\!\top}v\rangle+\|\delta\Psi_s^{\!\top}v\|^2\big]$,
which is positive whenever $\|\delta\Psi_s^{\!\top}v\|>2\|\Psi^{\!\top}v\|=2\sqrt{n\lambda_v}$.
\end{proposition}

\begin{theorem}[Feature-sum identity and bracket at $T$]
\label{thm:bracket}
$\Psi_T=\Psi_0+\sum_s\delta\Psi_s$, so for every unit $v$, with
$D_v=\sum_s\delta\Psi_s^{\!\top}v$ and $\ell_v=\|\Psi_0^{\!\top}v\|$,
\[
\tfrac1n\big(\|D_v\|-\ell_v\big)_+^2\;\le\;v^{\!\top}K_Tv\;\le\;\tfrac1n\big(\|D_v\|+\ell_v\big)^2 .
\]
$D_v$ splits exactly into coherent and propagation parts by
Proposition~\ref{prop:chain-identities}(d). On the directions the residual has
driven, $\ell_v$ is negligible and the lower bound is the value: measured
$0.90$--$1.00$ of $v^{\!\top}K_Tv$ for every rank $\ge60$ on the reference
problem [n=200]. Where $J$ is $C^1$ along the path,
$|\lambda_k(K_T)-\lambda_k(K_0)|\le\int_0^T C(s)\|E(s)\|\,ds$ with
$C=\frac2{n^2}\|\Psi\|_2^2\|\Psi'\|_2$ (Weyl), so every lifted eigenvalue at $T$
is bracketed between two functionals of the residual history.
\end{theorem}

\begin{theorem}[Stage error and the chain]
\label{thm:chain}
Where $J$ is $C^2$ along the path,
$\delta\Psi_s=\alpha_s\dot\Psi(t_{s-1})+R_s$ with
$\|R_s\|\le\frac{\alpha_s^2}2\sup\|\ddot\Psi\|$; choosing
$\alpha_s\|\ddot\Psi\|\le\varepsilon\|\dot\Psi\|$ makes each stage's first-order
increment accurate to relative $\varepsilon$, and the composition of the
stages is the exact trajectory. The quantity $\|\ddot\Psi\|/\|\dot\Psi\|$ is the
inverse of the stage length over which the linearisation at the current state
is trustworthy; its norm along the descent direction is $H_2$.
\end{theorem}
\begin{remark}
A numerical realisation of the chain (one eigendecomposition of the
$N\times N$ kernel per stage, the maps of
Proposition~\ref{prop:chain-identities}(c,d) applied to the network's own
units) reproduces the real gradient-flow run to four digits in $\|E\|$ and to
$0.1\%$ in every kernel eigenvalue at $t=8$, tracks it to $1$--$3\%$ at $T$ on
every design tested including depth~8 and width~512, and is converged to
three digits across a tenfold range of stage budgets [n=200]. It is gradient
flow on the sample integrated in large stages; its only use in this section is
as the exact reference against which every reduced statement is checked.
\end{remark}

\subsection{What a stage does to the weights: the spike theorem}
\label{ssec:chain-spikes}

For a hidden block whose outgoing weights are $W_{\ell+1}$, write
$S[i,q]=\sigma'(z_q(x_i))$ for the local derivative of the map following the
block, $M[i,q]$ for the backprop coefficient contracted with the drive (for the
last hidden block $M=U\,W_{\mathrm{out}}$, of rank $\le n_y$), and $A$ for the
block's inputs on the data. By Proposition~\ref{prop:chain-identities}(c),
$\Delta W=-\frac1n(S\odot M)^{\!\top}A$.

\begin{theorem}[Core and bulk]
\label{thm:spike}
\begin{enumerate}
\item[(a)] (\emph{at initialisation}) Split $S=h_0+(S-h_0)$ with $h_0$ the mean of
$\sigma'$ over the data. Then $\Delta W=\Delta W^{\mathrm{core}}+\Delta W^{\mathrm{bulk}}$ with
$\Delta W^{\mathrm{core}}=-\frac{h_0}nM^{\!\top}A$ of rank $\le n_y$, its left
singular space spanned by the rows of $W_{\mathrm{out}}$ and its right singular
space by $A^{\!\top}U$. If the block's input weights are independent of
$W_{\ell+1}$ (true at $\theta_0$), the expected bulk-to-core energy ratio is
$\rho=\mathrm{Var}(\sigma'(z))/\mathbb E[\sigma'(z)]^2$ at the preactivation scale.
The output block has rank $\le n_y$ and the input block rank $\le d_{\mathrm{in}}+1$ exactly.
\item[(b)] (\emph{any time, no independence}) For any rank-$r$ approximation
$S_r$ of $S$ with remainder $\tilde S$,
$\mathrm{rank}\big((S_r\odot M)^{\!\top}A\big)\le n_y r$ and
$\|(\tilde S\odot M)^{\!\top}A\|_F\le\|\tilde S\|_{\max}\|M\|_F\|A\|_2/n$. Since
$Z=A\,W_\ell^{\!\top}+b$, the numerical rank of $S=\sigma'(Z)$ is controlled by the
spectrum of the block's input Gram and by the Hermite coefficients of
$\sigma'$ at the current preactivation scale, both quantities the chain
carries; for deeper blocks the bound composes multiplicatively down the stack.
\end{enumerate}
\end{theorem}
\begin{proof}
(a) Expand $\sigma'$; the core is $h_0M$; the bulk's expected energy factorises
by independence of $S_{iq}$ and $M_{iq}$. (b) $\mathrm{rank}(X\odot Y)\le
\mathrm{rank}X\cdot\mathrm{rank}Y$, the entrywise Frobenius bound, the
column-space inclusion for $Z$, and the Hermite expansion.
\end{proof}

\begin{remark}[What was measured, n=200]
(a) is verified to within $10\%$ at initialisation for tanh, ReLU, GELU and
softplus; the bulk is $2$--$5\%$ of the core for tanh at unit gain and
$30$--$50\%$ for ReLU (whose local derivative is a $0/1$ pattern), yet $98$--$100\%$
of the first-order change lies in three singular directions for every
activation. Along the path the independence of (a) fails within a few steps,
as it must once the units adapt to the residual, while the rank persists: at
$t=0.3$ the per-unit-mean split gives a bulk ratio of $0.9$ against the
dispersion prediction of $0.05$, and the cumulative change is still $98\%$ in
three directions; the correlation-free bound (b) holds with $r=1$--$2$, i.e.\
the update lies $93\%$ in $3$ and $99.9\%$ in $9$ directions. Cores of
successive stages add coherently while the drive keeps its direction, which
is the cross-layer alignment measured earlier (randomising one layer's
rotation destroys the lift); the bulk-plus-spike spectrum of trained weights
is the random initialisation plus this coherent low-rank sum.
\end{remark}

\subsection{The richness expansion, its radius, and the chain as its resummation}
\label{ssec:chain-richness}

Let $\epsilon$ be any parameter that multiplies the weight velocity at fixed
residual dynamics in a suitable time: the inverse output scale $1/s$ with
$\tau=s^2t$ for centred outputs $\hat Y=s(N_\theta-N_{\theta_0})$; the Euler step
$\Delta$ of \S\ref{ssec:chain-apps}; a $\mu$P-type multiplier; an inverse-width
factor. Then
\[
\frac{dE}{d\tau}=-K(\theta)E,\qquad\frac{d\theta}{d\tau}=-\frac{\epsilon}{n}J(\theta)^{\!\top}E .
\]

\begin{theorem}[All-orders richness expansion]
\label{thm:richness}
For analytic parameter-free maps and each fixed $\tau$ there is $R(\tau)>0$ such
that every analytic observable of the state is a convergent power series in
$\epsilon$ for $|\epsilon|<R(\tau)$; in particular
\[
\frac{v^{\!\top}K(\tau)v}{v^{\!\top}K(0)v}-1=\sum_{k\ge1}a_k(\tau)\,\epsilon^k ,
\]
with coefficients determined recursively from $\theta_0$ and the lazy
trajectory alone: $a_0$ is kernel regression with $K(\theta_0)$,
$\theta^{(1)}(\tau)=-\frac1nJ_0^{\!\top}U_{\mathrm{lazy}}(\tau)$,
$a_1=2\langle\Psi_0^{\!\top}v,J'[\theta^{(1)}]^{\!\top}v\rangle/(n\lambda_v)$, and $a_k$
involves derivatives of $J$ at $\theta_0$ up to order $k$ along lower-order
displacements and the lower-order residual corrections (linear inhomogeneous
equations, variation of constants in the eigenbasis of $K_0$). No quantity along
a trained path enters at any order.
\end{theorem}

\begin{theorem}[Radius, and the chain as resummation]
\label{thm:radius}
$R(\tau)$ is bounded below by the stage criterion of Theorem~\ref{thm:chain}
applied at $\theta_0$, and along a direction $v$ the series is valid while
$\|\delta\Psi^{\!\top}v\|<\|\Psi_0^{\!\top}v\|$ (Proposition~\ref{prop:stage}). Where
this fails the lift along $v$ is of order one and no finite or infinite number
of terms of the series about $\theta_0$ represents it. Re-expanding at each
stage boundary (Theorem~\ref{thm:chain}) continues the series past $R(\tau)$
by moving the expansion point along the path; within $R(\tau)$ the two agree,
beyond it only the chain is defined.
\end{theorem}
\begin{proof}
Analytic dependence of ODE solutions on parameters; the recursion is the Taylor
expansion of the vector field in $\epsilon$; the radius statement is
Proposition~20 at the expansion point $\theta_0$ together with
Proposition~\ref{prop:stage}.
\end{proof}

\begin{remark}[Verification, n=200]
With $\epsilon=1/s$ and the coefficients through second order (including the
second-order weight change and the feedback of the lifted kernel on the lazy
residual), the predicted lift of init eigen-directions agrees with the
measured one to $0$--$16\%$ at output scales $s=2,4,8,16$ for as long as the
measured lift is below order one, including the early \emph{decrease} along
deep-tail directions and its sign change; the truncated series fails sharply
once the lift exceeds one, where Theorem~\ref{thm:radius} says it must. The
leading coefficient $a_1$ changes sign along $\tau$, so no universal power of
$s$ describes the delay to a given lift; the function is predicted, the
exponent does not exist.
\end{remark}

\subsection{Applications}
\label{ssec:chain-apps}

\paragraph{Euler versus direct.} Let $y_D=\Phi_\Delta(x)$ and
$y_E=(\Phi_\Delta(x)-x)/\Delta$ be the direct and the Euler target, trained from
the same $\theta_0$ with centred outputs, and write the direct residual on the
increment as $E_D=\Delta\tilde E$.
\begin{theorem}
\label{thm:euler}
$\tilde E$ and $E_E$ obey the same residual equation with their own kernels,
start from the same value, and the weight velocities satisfy
$\dot\theta_D=-\frac\Delta nJ^{\!\top}\tilde E$ against
$\dot\theta_E=-\frac1nJ^{\!\top}E_E$: the direct problem is the Euler problem with
the richness parameter multiplied by $\Delta$ (Theorem~\ref{thm:richness}
evaluated at a smaller argument). Hence (i) direct's stage lifts on the
increment are $\Delta^2$ times Euler's at equal normalised residual; (ii) as
$\Delta\to0$ at fixed budget the direct network is in the lazy regime on the
increment, its normalised error converges to the frozen-kernel error of the
increment target, and the held-rough bound of Section~\ref{sec:held-rough}
becomes tight for it; (iii) Euler's error is the full feature-learning error,
below that bound by the feature-learning gain, which grows with the budget
and with the increment's roughness. The direct network additionally spends a
$\Delta$-independent early drive on the identity component of its target.
\end{theorem}
The static argument of the previous sections (``the increment is a smoother
target'') is the frozen-kernel shadow of this: dividing by $\Delta$ multiplies the
drive by $1/\Delta$ and the lifts by $1/\Delta^2$, so Euler trains rich on the
increment while direct trains lazy on it. Statement (ii) is exact; (iii) is
conditional on the rate comparison below.

\paragraph{Progressive sharpening, the initial reduction phase, and the edge of stability.}
By Proposition~\ref{prop:stage},
$\dot\lambda_1=v_1^{\!\top}(\mathrm{Cross}+Q)v_1$ per stage with $Q\succeq0$. Along
the residual the cross term at $t=0$ equals $-\frac{2\alpha}{n^2}\langle
E,N''_{\theta}[g,g]\rangle$ with $g=J^{\!\top}E$, exactly, and it is negative on
$90$--$91\%$ of the eigen-directions for the raw residual $E=\hat Y_0-Y$ and on
$48$--$55\%$ (no systematic sign) for the centred residual $-Y$ [n=200]: the
initial sharpness-\emph{reduction} phase reported in the warmup literature is
the network unlearning its own initial function and is absent for centred
initialisations. Once learning has accumulated, the coherent part of the cross
term is the inner product of the accumulated core above with the stage core
below through the shared drive, positive while the drive keeps its direction
(measured positive on $80$--$100\%$ of modes from $t\approx1$ on): progressive
sharpening, faster with depth by composition ($5\times,11\times,20\times,86\times$
top gain at depths $2,3,6,6$ [n=200]). Taking GD steps as stages, above the
stability threshold the top component of the residual alternates in sign, the
stage cores along it alternate, their cross terms with the accumulated core
are negative at the core level, and $\lambda_1$ decreases until the step is
stable again: sharpening and self-stabilisation are the two signs of one term.
What is not proved is the sign of the init-feature part of the cross term
after the transient (Lemma~\ref{lem:sign} below).

\paragraph{Adam.} For $\dot\theta=-P(t)\nabla L$ with $P$ diagonal positive, (a)
of Proposition~\ref{prop:chain-identities} holds with $K_P=\Psi P\Psi^{\!\top}/n$,
(c) with the outer products multiplied entrywise by $P_\ell$, and
Proposition~\ref{prop:stage} and Theorem~\ref{thm:bracket} verbatim with
$P^{1/2}$-weighted features. The spike statement needs one more fact: $P$ is
built from the gradient block itself, so $P\odot G$ is a normalised sign
pattern of a rank-$\le n_y$ matrix, and the same mean-plus-fluctuation
argument as Theorem~\ref{thm:spike} gives its core fraction as the squared first
Hermite coefficient of the normaliser under the entry distribution of $G$:
$2/\pi$ for Gaussian-like entries, $1$ for coordinate-aligned (sparse) singular
vectors, with the time-averaged second moment lowering it further (measured
$0.68$--$0.81$, $0.95$--$1.00$, $0.52$ respectively on synthetic blocks). This
is the rotation-dependent property that distinguishes Adam from SGD.

\paragraph{Validation.} With $J_{\mathrm{val}}$ the Jacobian on validation
points, $\dot{\hat Y}_{\mathrm{val}}=-K_{vt}E$ with
$K_{vt}=J_{\mathrm{val}}J^{\!\top}/n$, exactly, so
$-\frac d{dt}\frac12\|E_{\mathrm{val}}\|^2=E_{\mathrm{val}}^{\!\top}K_{vt}(0)E+
E_{\mathrm{val}}^{\!\top}[K_{vt}(t)-K_{vt}(0)]E$: a lazy-transfer part and a
feature-learning-transfer part. The coherent part of the lift on validation
points is the training stage core transported by the cross activation Gram;
under equal unit statistics on the two samples the whole chain applies with
that substitution. Measured [n=200]: early both parts positive and the
validation rate $70$--$110\%$ of the training rate; late the lazy part
negative (the frozen kernel maps a residual living in sample-specific rough
modes into directions that increase validation error, the static held-rough
mechanism seen as a sign) while the lifted part stays positive. At the sample
sizes of interest, where training and validation agree, the sign pattern is a
prediction to be checked rather than a result.

\paragraph{Delayed feature learning (``grokking'').} The mechanism is
Theorems~\ref{thm:spike}--\ref{thm:bracket}: the residual drives a positive
lift along the directions it has visited, so a direction too slow for the
initial kernel to fit within the budget is fit later through the lift. The
onset is Theorem~\ref{thm:richness} with $\epsilon=1/s$: an init-computable
series, verified to $0$--$16\%$ in its range, with a proved radius; the delay
to a small lift is its inversion, and because $a_1$ changes sign no universal
exponent in $s$ exists, which accounts for the disagreement among reported
scalings. The event itself, a lift of order one or more along a misaligned
direction, lies outside $R(\tau)$ by construction and is a chain computation;
by \S\ref{ssec:chain-impossible} it is not predictable from the
initialisation in any other way. Rigorous timings in the literature (weight
decay, modular arithmetic, two-layer mean field) are special cases evaluated
inside their own radii.

\paragraph{The rate comparison, and the one lemma it rests on.}
\begin{lemma}[Sign of the cross term; status: derived for the core part, measured for the rest]
\label{lem:sign}
$v^{\!\top}\mathrm{Cross}_sv=\frac2n\big[\langle\Psi_0^{\!\top}v,\delta\Psi_s^{\!\top}v\rangle+\sum_{r<s}\langle\delta\Psi_r^{\!\top}v,\delta\Psi_s^{\!\top}v\rangle\big]$.
The second term's core-core part is nonnegative while the integrated residual
keeps its direction (shared drive); the first term is measured negative at
$t=0$ for raw residuals and sign-free for centred ones, positive after the
transient, and is not derived.
\end{lemma}
If the lemma holds from $t_0$ on, then $\dot L(t)\le-\langle E,K(t_0)E\rangle$ for
$t\ge t_0$: feature learning never slows the descent below what the kernel at
$t_0$ would produce on the current residual. A closed-form bound on
$\|E(T)\|$ needs in addition a persistence condition on the residual
direction, because $K(t)\succeq K(t_0)$ does not control mode-wise decay when
the eigenvectors rotate, and they rotate by about half in the relevant band by
mid-training [n=200].

\subsection{Why the late stage is not predictable from the initialisation}
\label{ssec:chain-impossible}

\begin{proposition}[Determinism; no chaos]
\label{prop:nochaos}
For analytic $N_\theta$ the gradient-flow trajectory is real-analytic in $t$, so
the history on any interval determines $E(T)$. The loss is a strict Lyapunov
function, so there are no recurrent sets and no sustained exponential
separation of nearby trajectories; what remains is finite-time stretching
along transiently negative-curvature directions and algebraic amplification
through the lift. [n=200: a $10^{-7}$ relative perturbation of $\theta_0$ grows by
a factor $1.8$ (random direction), $4.6$ (tail-feature direction) and shrinks to
$0.65$ (top-feature direction) by $t=76$; continuation to $T$ in progress.]
\end{proposition}

So the obstruction is conditioning, and it has three parts, each a theorem.

\begin{theorem}[Hidden state]
\label{thm:hidden}
Let $A$ and $\Delta$ be the input activations and the outgoing backprop
signals of a hidden-hidden block on the data at $\theta_0$. For every
$N=UV^{\!\top}$ with $U^{\!\top}\Delta=0$ and $V^{\!\top}A=0$ (such $N$ exist whenever the
widths exceed the ranks, e.g.\ $H>n\,n_y$, and have any norm), the
initialisation $\theta_0+N$ has the same unit states, output, residual, loss,
gradient and kernel on the data, and differs only in the propagation part of
Proposition~\ref{prop:chain-identities}(d). Hence $E(T)$ is not a function of
$(K_0,E_0,Y)$, nor of any finite number of their time derivatives at $0$ when
the widths allow the corresponding null spaces. [Verified: with
$\|N\|_F=\|W\|_F$ the states agree to $5\cdot10^{-15}$ and the kernel increments
after one step differ by $1\%$.]
\end{theorem}

\begin{proposition}[Visibility]
\label{prop:visibility}
Let $B$ be the modes of $K(t_0)$ carrying a fraction $\phi$ of $\|E(T)\|^2$. Their
share of the loss-decrease rate at $t_0$ is at most
$\max_B\lambda_k(t_0)/\lambda_{\mathrm{eff}}(t_0)$, and in the fixed-basis
solution their components are constant on $[0,t_0]$ to relative order
$\lambda_k(t_0)t_0$. [n=200: median $1.4\cdot10^{-6}$ at $t_0=0.01T$ and
$10^{-2}$ at $t_0=0.1T$; share of the rate $8\cdot10^{-7}$ and $7\cdot10^{-4}$.]
\end{proposition}

\begin{proposition}[Amplification]
\label{prop:amplification}
The same modes are later lifted by the gain of
Theorem~\ref{thm:bracket} (measured $10^2$--$10^4$), coherently across blocks,
quadratically in the accumulated drive, and $E_k(T)\sim\exp(-\!\int\lambda_k)$ is
exponentially sensitive to the lifted eigenvalue. In the fixed-basis model the
sensitivity of $\log|E_k(T)|$ to a mode's growth exponent is
$2\Lambda_k(T)\log(T/t_0)$ with $\Lambda_k(T)\approx1$ for the modes that
matter.
\end{proposition}

\begin{corollary}[The conditioning statement]
\label{cor:conditioning}
To predict $E(T)$ to relative precision $\varepsilon$ from the history up to
$t_0$ the tail components must be known to relative precision of order
$\varepsilon\,\lambda_k(t_0)t_0$ divided by the amplification, i.e.\ about
$10^{-7}$ at $t_0=0.01T$ for $\varepsilon=0.1$ [n=200]: five orders of magnitude
below the initialisation-to-initialisation variation of the early history
($2\cdot10^{-2}$) and below single precision. The late state is a smooth,
deterministic function of the early state whose relevant derivatives are
enormous exactly in the directions that are unobservable early.
\end{corollary}

The three parts have a common root: the drive is the residual, the residual
lives in the slow modes of the current kernel, and the slow modes are both the
ones that have not moved yet (invisible) and the ones feature learning
targets (amplified). Feature learning is the mechanism by which the network
amplifies precisely the information that no early observable contains, and
that information is seeded by the initialisation through
Theorem~\ref{thm:hidden}. This is why every early predictor we tried (frozen
kernel, gain-law extrapolation, residual front, the loss itself) orders
designs within $1.1\times$ at chance and within $1.3\times$ at $65$--$76\%$, why
seed pairs that agree to $2\%$ for the first tenth of training end $15\%$
apart, and why the only object that reproduces the late error is the
trajectory carried forward [n=200].

\subsection{What a statistical description can promise}
\label{ssec:chain-statistical}

The impossibility above is a statement about one network. A statistical
description does not predict one network; it predicts the distribution over
initialisations, and by Theorem~\ref{thm:hidden} that is the correct object,
because the quantity that defeats prediction is the draw of the hidden state.
Three consequences follow, two of them proved and one to be measured.

\emph{The irreducible uncertainty of any statistical prediction for a single
run is the spread of the late error over initialisations}, $15\%$ median here
[n=200]. A design comparison finer than that spread is not well posed for one
run of each design, whatever the method; it is well posed for the expected
error, which is a statistic.

\emph{The coherent chain is the statistical approximation}: replace the
network's own units by draws from the unit distribution and the propagation
part by its average, and Propositions~\ref{prop:chain-identities}(b,c) and
Theorem~\ref{thm:spike}(a) close on the data. On the one case run to
convergence [H=96, depth 2, n=200] it reproduced the structure, the gain law
and the ordering of designs, with a $2\times$ bias in the late error and a
$4$--$5\times$ excess in the tail gain traced to the averaged propagation term and
to the transfer of the ensemble's lift; at depth~3 the bias was larger
($5\times$ in the tail gain). It was not run at larger width. Whether the
bias and the seed spread decrease with width, as the propagation term being
a finite-width fluctuation suggests, or persist with depth, as the depth-3
result warns, is the single most consequential unmeasured quantity in this
work, and it decides the question in the next paragraph.

\emph{The diagnostic that is well posed is the expected late-stage ability of an
architecture.} If the spread over initialisations shrinks with width and the
coherent chain converges to the exact one in that limit, then ranking
architectures by their expected late error is both meaningful and computable
without the particular initialisation, and the finite-width exponent of the
convergence is the error bar. Everything in this section is set up for that
measurement: the exact chain for the per-initialisation truth, the coherent
chain for the expectation, and the seed spread for the floor. The ranker
built from the identities (lazy error, $H_2$, top eigenvalue, output scale)
reached a $1.3\times$ resolution out of sample [n=200], which is the expected
accuracy of a one-Jacobian statistic and should be read as a lower bound on
what the expected-error ranking can do.

\subsection{Scope and what to measure}
\label{ssec:chain-scope}

Every number in this section comes from $n=200$ training points on the
Lorenz-63 one-step map, at which the training set is nearly fit within the
budget, the residual-adapted late phase is short, the depth-$6$--$8$ ``floor'' is a
sample-size artefact, and the negative lazy-transfer sign in the validation
split is overfitting of a $200$-point sample. The theorems carry $n$ as a symbol
and are unaffected; the verifications, bulk ratios, exponents, signs, transfer
fractions and the ranker's fits are all $n=200$ statements and are to be redone
at a sample size that leaves a genuine late stage before any is quoted. The
exact chain does not scale to large $n$ (one $N\times N$ eigendecomposition per
stage); the ground-truth runs do, so the redo is training-only, and it should
include, at that $n$: the seed spread of the late error against width and
depth; the gap between the own-initialisation and the other-initialisation
chain against width and depth; the cross-term sign and the spike bound along
the trajectory; the validation split; and the Euler-versus-direct comparison
at several $\Delta$ from a shared initialisation.
